%------------------------------------------------------------------------------%

\usepackage{apresentacao}
\usepackage{tabto}

\graphicspath{{./figs/}}

%------------------------------------------------------------------------------%
\newcommand{\mytitle}   {Séries de Taylor}
\newcommand{\mysubtitle}{Definição e propriedades}
\title   {\mytitle}
\subtitle{\mysubtitle}
\author  {\large Luis Alberto D'Afonseca}
\date    {}

%------------------------------------------------------------------------------%
\begin{document}

%------------------------------------------------------------------------------%
\begin{frame}
    \titlepage
\end{frame}

%------------------------------------------------------------------------------%
\section{Séries de Taylor}

%------------------------------------------------------------------------------%
\begin{frame}{Séries de Potência e Séries de Taylor}
  \begin{itemize}[<+->]
    \item \emph{Séries de Potências} \\[3mm]
          Dada uma Série de Potências construímos uma função \\[11mm]
    \item \emph{Séries de Taylor}   \\[3mm]
          Dada uma função construímos uma Série de Potências
  \end{itemize}
\end{frame}

%------------------------------------------------------------------------------%
\begin{frame}{Séries de Taylor e Séries de Maclaurin}
  \begin{itemize}[<+->]
    \item \emph{Séries de Taylor}   \\[3mm]
    Centrada em um ponto \emph{$a$ qualquer} \\[11mm]
    \item \emph{Séries de Maclaurin}   \\[3mm]
    Série de Taylor centrada na origem, \emph{$a=0$}
  \end{itemize}
\end{frame}

%------------------------------------------------------------------------------%
\section{Coeficientes da Série de Taylor}

%------------------------------------------------------------------------------%
\begin{frame}{Construção da Série de Taylor}
  Supomos que uma função pode ser escrita como uma Série de Potências
  \[
    f(x) = \sum_{n=0}^{\infty} c_n (x-a)^n
  \]
\end{frame}

%------------------------------------------------------------------------------%
\begin{frame}[t]{Construção da Série de Taylor}
  \vspace{\baselineskip}
  \begin{align*}
    \action<+->{f(x)                 & = \sum_{n=0}^{\infty} c_n (x-a)^n      \\[2mm]}
    \action<+->{\phantom{f^{(3)}(x)} & = \memph{c_0} + c_1(x-a) + c_2(x-a)^2 + \cdots \\[5mm]}
    \action<+->{f(a)                 & = \memph{c_0}}
  \end{align*}
  \action<+->{
  \[ \memph{c_0} = f(a) = \frac{f(a)}{0!} \]}
\end{frame}

%------------------------------------------------------------------------------%
\begin{frame}[t]{Construção da Série de Taylor}
  \vspace{\baselineskip}
  \begin{align*}
    \action<+->{f'(x)
      & \vphantom{\sum_{n=0}^{\infty}}
        = \frac{d}{dx}\Big[c_0
        + \memph{c_1}(x-a)
        + c_2(x-a)^2
        + \cdots\Big] \\[2mm]}
    \action<+->{\phantom{f^{(3)}(x)}
      & = \memph{c_1}
        + 2c_2(x-a)
        + 3c_3(x-a)^2
        + \cdots \\[5mm]}
    \action<+->{f'(a) & = \memph{c_1}}
  \end{align*}
  \action<+->{
  \[
    \memph{c_1} = f'(a) = \frac{f'(a)}{1!}
  \]}
\end{frame}

%------------------------------------------------------------------------------%
\begin{frame}[t]{Construção da Série de Taylor}
  \vspace{\baselineskip}
  \begin{align*}
    \action<+->{f''(x)
      & \vphantom{\sum_{n=0}^{\infty}}
        = \frac{d}{dx}\Big[ c_1
        + 2\memph{c_2}(x-a)
        + 3c_3(x-a)^2
        + \cdots\Big] \\[2mm]}
    \action<+->{\phantom{f^{(3)}(x)}
      & = (2\times 1)\memph{c_2}
        + (3\times 2)c_3(x-a)
        + (4\times 3)c_4(x-a)^2
        + \cdots \\[5mm]}
    \action<+->{f''(a)
      & = (2\times 1)\memph{c_2}}
  \end{align*}
  \action<+->{\[
    \memph{c_2} = \frac{f''(a)}{2!}
  \]}
\end{frame}

%------------------------------------------------------------------------------%
\begin{frame}[t]{Construção da Série de Taylor}
  \vspace{\baselineskip}
  \begin{align*}
    \action<+->{f^{(3)}(x)
      & \vphantom{\sum_{n=0}^{\infty}}
        = \frac{d}{dx}\Big[(2\times 1)c_2
        + (3\times 2)\memph{c_3}(x-a)
        + (4\times 3)c_4(x-a)^2
        + \cdots\Big]
      \\[2mm]}
    \action<+->{\phantom{f^{(3)}(x)}
      & = (3\times 2\times 1)\memph{c_3}
        + (4\times 3\times 2)c_4(x-a)
        + (5\times 4\times 3)c_5(x-a)^2
        + \cdots \\[5mm]}
    \action<+->{f^{(3)}(a)
      & = (3\times 2\times 1)\memph{c_3}}
  \end{align*}
  \action<+->{
  \[
    \memph{c_3} = \frac{f^{(3)}(a)}{3!}
  \]}
\end{frame}

%------------------------------------------------------------------------------%
\begin{frame}[t]{Construção da Série de Taylor}
  \vspace{\baselineskip}
  \begin{align*}
    \action<+->{f^{(4)}(x)
      &  \vphantom{\sum_{n=0}^{\infty}}
        = \frac{d}{dx}\Big[(3\times 2\times 1)c_3
        + (4\times 3\times 2)\memph{c_4}(x-a)
        + (5\times 4\times 3)c_5(x-a)^2
        + \cdots\Big]
       \\[2mm]}
    \action<+->{\phantom{f^{(3)}(x)}
      & = (4\times 3\times 2\times 1)\memph{c_4}
        + (5\times 4\times 3\times 2)c_5(x-a)
        + \cdots \\[5mm]}
    \action<+->{f^{(4)}(a)
      & = (4\times 3\times 2\times 1)\memph{c_4}}
  \end{align*}
  \action<+->{
  \[
    \memph{c_4} = \frac{f^{(4)}(a)}{4!}
  \]}
\end{frame}

%------------------------------------------------------------------------------%
\begin{frame}{Construção da Série de Taylor}
  Generalizando
  \[
    f^{(k)}(a) = k!c_k  \qquad\qquad c_k = \frac{\;f^{(k)}(a)\;}{k!}
  \]

  \pause
  \vspace{\baselineskip}
  Se    \(\quad f(x) = \sum_{n=0}^{\infty} c_n (x-a)^n \quad \)
  então \(\quad c_k = \frac{\;f^{(k)}(a)\;}{k!} \quad \)
\end{frame}

%------------------------------------------------------------------------------%
\section{Definição da Série de Taylor}

%------------------------------------------------------------------------------%
\begin{frame}{Série de Taylor}
  Se $f$ é uma função infinitamente derivável \\
  sua \emph{Série de Taylor} centrada em $\;x=a\;$ é
  \begin{multline*}
    \sum_{n=0}^\infty \frac{\;f^{(n)\;}(a)}{n!}\,(x-a)^n = \\
    f(a)\;+\;f'(a)(x-a)\;+\;\frac{f''(a)}{2!}(x-a)^2\;+\cdots+\;\frac{f^{(n)}(a)}{n!}(x-a)^n\;+\cdots
  \end{multline*}

  \pause
  \vspace{\baselineskip}
  Observe que \emph{não} escrevemos que $f(x)$ é igual a sua série de Taylor
\end{frame}

%------------------------------------------------------------------------------%
\begin{frame}{Série de Maclaurin}
  Uma Série de Maclaurin é uma Série de Taylor centrada em zero, $a=0$
  \[
    \sum_{n=0}^\infty \frac{\;f^{(n)\;}(0)}{n!}\,x^n
  \]
\end{frame}

%------------------------------------------------------------------------------%
\begin{frame}{Polinômios de Taylor}
  Se $f$ é uma função com $k$ derivadas \\
  seu \emph{Polinômio de Taylor} de ordem $k$ centrado em $\;x=a\;$ é
  \[
    \sum_{n=0}^k \frac{\;f^{(n)\;}(a)}{n!}\,(x-a)^n
      \;=\; f(a)
      \;+\;f'(a)(x-a)
      \;+\cdots+\; \frac{f^{(k)}(a)}{k!}(x-a)^k
    \]
\end{frame}

%------------------------------------------------------------------------------%
\section{Exemplo}

%------------------------------------------------------------------------------%
\begin{frame}{Exemplo}
  Calcular a Série de Taylor da função \emph{$f(x)=e^x$} centrada em \emph{$a=0$}

  \pause
  \vspace{0.5\baselineskip}
  Derivando
  \[
    f^{(k)}(x) = \frac{d^k}{dx^k}\,e^x = e^x
  \]
  \pause
  Avaliando em $a=0$
  \[
    f^{(k)}(0) = e^0 = 1
  \]
\end{frame}

%------------------------------------------------------------------------------%
\begin{frame}{Exemplo -- Calculando os Coeficientes}
  \begin{align*}
    \action<+->{c_0 & = \frac{f^{(0)}(0)}{0!} = 1  \\[3mm]}
    \action<+->{c_1 & = \frac{f^{(1)}(0)}{1!} = 1 \\[3mm]}
    \action<+->{c_2 & = \frac{f^{(2)}(0)}{2!} = \frac{1}{2} \\[7mm]}
    \action<+->{c_k & = \frac{f^{(k)}(0)}{k!} = \frac{1}{k!}}
  \end{align*}
\end{frame}

%------------------------------------------------------------------------------%
\begin{frame}{Exemplo -- Construindo a Série}
  \begin{align*}
    \action<+->{S(x) & = \sum_{n=0}^\infty \frac{f^{(n)}(a)}{n!}(x-a)^n \\[3mm]}
    \action<+->{     & = \sum_{n=0}^\infty \frac{f^{(n)}(0)}{n!}x^n     \\[3mm]}
    \action<+->{     & = \sum_{n=0}^\infty \frac{x^n}{n!}}
    \action<+->{     \;=\; 1 + x + \frac{x^2}{2} + \frac{x^3}{3!} + \cdots + \frac{x^n}{n!}+\cdots}
  \end{align*}
\end{frame}

%------------------------------------------------------------------------------%
\begin{frame}{Exemplo -- Convergência}
\begin{itemize}[<+->]
  \item Vimos que a série converge para todo $x\in\R$ \\[5mm]
  \item A série converge para a função? \\[5mm]
  \[
     e^x \;\;{\color{red}\stackrel{?}{ = }}\;\; \sum_{n=0}^\infty \frac{\;x^n\;}{n!}
  \]
\end{itemize}
\end{frame}

%------------------------------------------------------------------------------%
\begin{frame}{\mytitle}{\mysubtitle}
  \tableofcontents
\end{frame}

%------------------------------------------------------------------------------%
\end{document}
%------------------------------------------------------------------------------%
